\documentclass[11pt]{article}
\usepackage{mathblatt}
% Präambel des Setzers (werkzeuge/setzer-praeambel.tex, wortgleich); dazu \gteil (Muster T2U)
\makeatletter
\setlength{\mbpfbe}{0mm}% keine Punkte (pruefung.md § 5: nur beim Prüfstein)
\newlength{\gEin}\setlength{\gEin}{0pt}
\let\gSchrift\relax
\renewcommand{\pfkreuzzeile}[1]{\par\noindent\hangindent3.2em\hangafter1\hspace*{1.6em}$\square$\ #1\par}
\newcommand{\gkreuz}[1]{$\square$\ #1\hspace{2em}}
\newcommand{\gkreuzzeile}[2]{\par\noindent\hangindent2.2em\hangafter1\mbox{#1\,$\square$}\ #2\par}% Buchstabe, Option: eine je Zeile (neben Skizze)
% links, was man ansieht, rechts, was man tut; Spaltengrenze überall an derselben Stelle (gemeinsam.md „Satz“)
% v0.17: Fußeintrag „Nr: Wert“ – kurz in einer Box (nie getrennt), lang umbrechbar mit der Nummer am Anfang
% (vorher \mbox: lange Lösungen liefen über den Rand, AS40 138 pt); robust, weil \fusshilfe in die .aux schreibt
\NewDocumentCommand{\gfuss}{m m}{\sbox0{#1:~#2}\ifdim\wd0<0.45\textwidth\mbox{#1:~#2}\else\mbox{#1:}\nobreak\ #2\fi}
\newsavebox{\gbox}\newlength{\gLinks}
% v0.10 (Lehrer 13.10.): keine starre Halbierung – die linke Spalte ist so breit wie die Figur (+ 6 mm), die
% Antwort steht direkt daneben; ist die Figur breiter als 0,62 Zeile, steht die Antwort darunter
\newcommand{\gzwei}[2]{\par\smallskip\noindent\sbox{\gbox}{#1}%
  \setlength{\gLinks}{\dimexpr\wd\gbox+6mm\relax}%
  \ifdim\gLinks>0.62\linewidth
    \usebox{\gbox}\par\smallskip\noindent\begin{minipage}{\linewidth}#2\end{minipage}%
  \else
    \begin{minipage}[t]{\gLinks}\vspace{0pt}\usebox{\gbox}\end{minipage}%
    \begin{minipage}[t]{\dimexpr\linewidth-\gLinks\relax}\vspace{0pt}\raggedright #2\end{minipage}%
  \fi\par}
% v0.9 (12.10., pruefung.md § 5 „Marke mit B, Bedingungen rechts“): drittes Argument von pfaufg = Bedingungen
% der Prüfung (v0.10: \pfohnetr = TR durchgestrichen, „FOR“), klein grau im rechten Rand auf der Grundlinie der Nummer, Zeilen mit \\
\newcommand{\gbedrand}[2]{\ifblank{#2}{}{\rlap{\hspace*{\dimexpr#1+2mm\relax}%
  \parbox[t]{15mm}{\raggedright\scriptsize\leavevmode\color{mbgrau}\strut#2}}}}
% Teil eines Bündels (a, b …): Marke links im Rand, Bedingungen rechts im Rand, auf der Zeile des Teils
\newcommand{\gteilrand}[2]{\llap{\mbpfrandmarke{#1}\hspace*{\mbpfnr}}\gbedrand{\linewidth}{#2}}
% Teil eines Bündels: Buchstabe in eigener Spalte, Inhalt eingerückt (\gEin = Einzug, damit die Spaltengrenze von
% \gzwei an derselben Stelle bleibt)
\newenvironment{gbteil}[3]{\par\smallskip\noindent\gteilrand{#2}{#3}\makebox[\mbpfteil][l]{\textbf{#1}}%
  \begin{minipage}[t]{\dimexpr\linewidth-\mbpfteil\relax}\setlength{\gEin}{\mbpfteil}\raggedright\strut\ignorespaces}%
 {\end{minipage}\par}
% Grafik höchstens so breit wie die Spalte (Bündel nebeneinander, Koordinatensystem ohne Karo)
\newsavebox{\gpbox}
\newcommand{\gpassend}[1]{\sbox{\gpbox}{#1}\ifdim\wd\gpbox>\linewidth\typeout{GPASSEND: verkleinert \the\wd\gpbox\space auf \the\linewidth}\resizebox{\linewidth}{!}{\usebox{\gpbox}}\else\usebox{\gpbox}\fi}
% Skizzenfeld (v0.10, Lehrer 13.10.): Karo ab dem linken Rand über die Zeile, „Skizze“ klein grau oben links im Feld
% (wie \pfskizzenfeld der Vorlage); der rechte Rand bleibt den Bedingungen
\newcommand{\gskizze}[1]{\par\smallskip\noindent\begin{tikzpicture}
  \clip (0,0) rectangle ({\linewidth-1mm},{#1*10mm});
  \draw[black!16,line width=0.3pt,step=5mm] (0,0) grid ({\linewidth},{#1*10mm});
  \node[anchor=north west,inner sep=2pt,font=\scriptsize,text=mbgrau,fill=white] at (0,{#1*10mm}) {Skizze};
  \end{tikzpicture}\par}
\RenewDocumentEnvironment{pfaufg}{m m m}{%
  \begin{lrbox}{\mbaufgabebox}\begin{minipage}[t]{\linewidth}%
  \gSchrift\noindent\hspace*{\gEin}\mbpfrandmarke{#1}\gbedrand{\linewidth-\gEin}{#3}%
  \begin{minipage}[t]{\mbpfnr}\strut\textbf{#2}\end{minipage}%
  \begin{minipage}[t]{\dimexpr\linewidth-\gEin-\mbpfnr\relax}\raggedright\strut\ignorespaces}%
 {\end{minipage}%
  \end{minipage}\end{lrbox}%
  \setlength{\mbaufgabehoehe}{\dimexpr\ht\mbaufgabebox+\dp\mbaufgabebox\relax}%
  \ifdefempty{\gkopfnach}{\par\addvspace{5pt}\Needspace*{\mbaufgabehoehe}}{%
    % Stufenkopf, der mit dieser (ersten) Aufgabe zusammenbleibt: Platz für beide verlangen, dann den Kopf setzen
    % (12.10., Quadratische: ein Kopf stand allein unten, die Aufgabe mit Koordinatensystem auf der nächsten Seite)
    % v0.17 (AS40 „Terme deuten“ allein unten): hinter dem Kopf steht das Sprungziel (\special, ein Knoten, nach
    % dem die Seite brechen durfte) – darum \nobreak nach dem Kopf; Vorhalt 56 pt = Kopf mit Linie und Abständen
    \par\Needspace*{\dimexpr\mbaufgabehoehe+56pt\relax}\gkopfnach\par\nobreak\gdef\gkopfnach{}\addvspace{5pt}}%
  \noindent\usebox{\mbaufgabebox}\par\addvspace{5pt}}
% Stufenkopf vormerken: \gkopfvor{<Kopf>} – die nächste pfaufg setzt ihn, sobald beide zusammen Platz haben
\def\gkopfnach{}
\newcommand{\gkopfvor}[1]{\gdef\gkopfnach{#1}}
% Teilaufgabe mit Skizze links in gzwei: Buchstabe in eigener Spalte, darunter Text und Skizze
\newcommand{\gteil}[2]{\makebox[\mbpfteil][l]{\textbf{#1}}\parbox[t]{\dimexpr0.5\textwidth-\gEin-\mbpfnr-5mm-\mbpfteil\relax}{\raggedright #2}}
% Stufenkopf der Vorlage (mathblatt.sty 2026-10-12: Linie, Abstand, 13 pt, \pfjahre) ohne ihren Vorhalt von
% 14 Zeilen: der Setzer fordert vor dem Kopf 0,16 Seitenhöhe (Kopf + erste Aufgabe), mehr nicht
\setlength{\mbpfstufevorhalt}{0pt}
\makeatother
\newcommand{\karo}[1]{\par\smallskip\noindent\begin{tikzpicture}
  \clip (0,0) rectangle ({\linewidth-1mm},{#1*10mm});
  \draw[black!16,line width=0.3pt,step=5mm] (0,0) grid ({\linewidth},{#1*10mm});
  \end{tikzpicture}\par}
\newcommand{\kk}[1]{$\square$\,#1\hspace{0.9em}}
\newcommand{\linie}[1]{\,{\color{black!45}\rule[-1pt]{#1}{0.4pt}}\,}
\newcommand{\pkt}{\textperiodcentered{}}
% v0.14 (14.10.): Stufenkopf mit fünf Punkten (Option = Jahre für die Bestellkarte) oder „zuletzt 19“
\NewDocumentCommand{\gpunkte}{o m}{#2}
\newcommand{\gzuletzt}[1]{zuletzt #1}
% v0.7 (12.10.): Schülerfrage unter dem Stufenkopf; Ankreuzoption als Bild (Buchstabe und Kästchen darunter)
\newcommand{\gfrage}[1]{\par\nobreak\noindent\hspace*{\mbpfnr}{\small\itshape #1}\par\nobreak\addvspace{2pt}}
\newlength{\gbildbreite}\setlength{\gbildbreite}{0.32\linewidth}% Setzer setzt sie je Reihe (vier Bilder: 0,24)
\newcommand{\gbildoption}[3]{\begin{minipage}[t][#3][b]{\gbildbreite}\centering #2\par\smallskip\mbox{#1\,$\square$}\end{minipage}}
% Lineare (12.10.): Bild mit Buchstabe ohne Kästchen (Diagramm zuordnen); Aussage mit „wahr / falsch“ zum Ankreuzen
\newcommand{\gbildname}[3]{\begin{minipage}[t][#3][b]{\gbildbreite}\centering #2\par\smallskip\mbox{\textbf{#1}}\end{minipage}}
\newcommand{\gwahrfalsch}[2]{\par\noindent\hangindent2.2em\hangafter1\makebox[2.2em][l]{(#1)}#2\hfill\mbox{wahr $\square$\quad falsch $\square$}\par}
\begin{document}

\pfheftstil
\pflaufkopf[Abitur GK \pkt{} Prüfung]{Extrempunkte \pkt{} Lösungen}{AA50}\pfseitevon
\pfheftkopf{Extrempunkte}{Lösungen}
\pfstufenzeile{Extremstelle nachweisen \pkt{} Extrempunkte gesucht \pkt{} Hochpunkt mit e-Funktion gesucht}

\Needspace*{0.16\textheight}\pfstufekopf{Extremstelle nachweisen}{\gpunkte[23 22]{\textcolor{black!25}{$\circ$}\,\textcolor{black!25}{$\circ$}\,\textcolor{black!25}{$\circ$}\,\textcolor{black!80}{$\bullet$}\,\textcolor{black!80}{$\bullet$}}}\special{pdf:dest (stufe-extremstelle-nachweisen) [@thispage /XYZ @xpos @ypos null]}
\begin{pfloesungb}{}
\lzb{1.}{}{h$''$ = $\tfrac{3}{4}$ t $-$ 18; Maximum; h(12) = 864 m}{h$'$(t) = $\tfrac{3}{8}$ t² $-$ 18t + 162; h$''$(t) = $\tfrac{3}{4}$ t $-$ 18; h$'$(12) = 54 $-$ 216 + 162 = 0 und h$''$(12) = $-$9 < 0 $\Rightarrow$ Maximum; h(12) = 216 $-$ 1296 + 1944 = 864 m}
\end{pfloesungb}
\begin{pfloesungb}{}
\lzb{2.}{}{waagerechte Tangente; f($-$2 + $\sqrt{13}$) $\approx$ $-$4,78}{($-$2 + $\sqrt{13}$)² = 17 $-$ 4$\sqrt{13}$ $\Rightarrow$ $\tfrac{1}{6}$ · (17 $-$ 4$\sqrt{13}$) + $\tfrac{2}{3}$ · ($-$2 + $\sqrt{13}$) $-$ $\tfrac{3}{2}$ = $\tfrac{17}{6}$ $-$ $\tfrac{8}{6}$ $-$ $\tfrac{9}{6}$ = 0 $\Rightarrow$ f'($-$2 + $\sqrt{13}$) = 0 $\Rightarrow$ waagerechte Tangente; f($-$2 + $\sqrt{13}$) = $\tfrac{1}{3}$ · (8 $-$ 4$\sqrt{13}$) · $e^{\tfrac{-2 + \sqrt{13}}{2}}$ = $\tfrac{4}{3}$ · (2 $-$ $\sqrt{13}$) · $e^{\tfrac{\sqrt{13}}{2} - 1}$ $\approx$ $-$4,78}
\end{pfloesungb}
\begin{pfloesungb}{}
\lzb{3.}{}{f$'$ = $-$2; Hochpunkt ($-$2|e²) $\approx$ ($-$2|7,39)}{$e^{-x}$ > 0 $\Rightarrow$ f$'$(x) = 0 nur für x = $-$2; f$'$($-$3) = $e^3$ > 0, f$'$($-$1) = $-$e < 0 $\Rightarrow$ Vorzeichenwechsel von + nach $-$ $\Rightarrow$ Hochpunkt ($-$2|e²) $\approx$ ($-$2|7,39)}
\end{pfloesungb}

\Needspace*{0.16\textheight}\pfstufekopf{Extrempunkte gesucht}{\gpunkte[25 24 23 22]{\textcolor{black!25}{$\circ$}\,\textcolor{black!80}{$\bullet$}\,\textcolor{black!80}{$\bullet$}\,\textcolor{black!80}{$\bullet$}\,\textcolor{black!80}{$\bullet$}}}\special{pdf:dest (stufe-extrempunkte-gesucht) [@thispage /XYZ @xpos @ypos null]}
\begin{pfloesungb}{}
\lzb{4.}{}{Extremstellen 0 und 4}{f$'$(x) = 3x² $-$ 12x = 3x · (x $-$ 4) = 0 $\Rightarrow$ Extremstellen 0 und 4}
\end{pfloesungb}
\begin{pfloesungb}{}
\lzb{5.}{}{x = 4; Tiefpunkt (0|0); Hochpunkt (4|$\tfrac{16}{9}$)}{f$'$(x) = $-\tfrac{1}{6}$ x² + $\tfrac{2}{3}$ x = $-\tfrac{1}{6}$ x · (x $-$ 4) = 0 $\Rightarrow$ x = 0 oder x = 4; f$''$(x) = $-\tfrac{1}{3}$ x + $\tfrac{2}{3}$: f$''$(0) = $\tfrac{2}{3}$ > 0 $\Rightarrow$ Tiefpunkt (0|0); f$''$(4) = $-\tfrac{2}{3}$ < 0 $\Rightarrow$ Hochpunkt (4|$\tfrac{16}{9}$)}
\end{pfloesungb}
\begin{pfloesungb}{}
\lzb{6.}{}{x = 4; Hochpunkt (2|4); Tiefpunkt (4|0)}{f$'$(x) = 3x² $-$ 18x + 24 = 3(x $-$ 2)(x $-$ 4) = 0 $\Rightarrow$ x = 2 oder x = 4; f$''$(x) = 6x $-$ 18: f$''$(2) = $-$6 < 0 $\Rightarrow$ Hochpunkt (2|4); f$''$(4) = 6 > 0 $\Rightarrow$ Tiefpunkt (4|0)}
\end{pfloesungb}
\begin{pfloesungb}{}
\lzb{7.}{}{x = $\pm\sqrt{3}$; Hochpunkt (0|2,5); Tiefpunkte ($-\sqrt{3}$|$-$2) und ($\sqrt{3}$|$-$2)}{f'(x) = 2x³ $-$ 6x = 2x(x² $-$ 3) = 0 $\Rightarrow$ x = 0 oder x = $\pm\sqrt{3}$; f''(x) = 6x² $-$ 6: f''(0) = $-$6 < 0 $\Rightarrow$ Hochpunkt (0|2,5); f''($\pm\sqrt{3}$) = 12 > 0 $\Rightarrow$ Tiefpunkte ($-\sqrt{3}$|$-$2) und ($\sqrt{3}$|$-$2)}
\end{pfloesungb}
\begin{pfloesungb}{}
\lzb{8.}{}{x = 3; Hochpunkt ($-$3|3), Tiefpunkt (3|$-$3); y = $-$x, Winkelhalbierende des II. und IV. Quadranten}{f$'$(x) = $\tfrac{1}{6}$ x² $-$ $\tfrac{3}{2}$ = 0 $\Rightarrow$ x² = 9 $\Rightarrow$ x = $-$3 oder x = 3; f($-$3) = 3, f(3) = $-$3; f$''$(x) = $\tfrac{1}{3}$ x $\Rightarrow$ Hochpunkt ($-$3|3), Tiefpunkt (3|$-$3); Steigung ($-$3 $-$ 3)/(3 + 3) = $-$1, durch den Ursprung $\Rightarrow$ y = $-$x, Winkelhalbierende des II. und IV. Quadranten}
\end{pfloesungb}

\Needspace*{0.16\textheight}\pfstufekopf{Hochpunkt mit e-Funktion gesucht}{\gpunkte[25 24]{\textcolor{black!25}{$\circ$}\,\textcolor{black!80}{$\bullet$}\,\textcolor{black!80}{$\bullet$}\,\textcolor{black!25}{$\circ$}\,\textcolor{black!25}{$\circ$}}}\special{pdf:dest (stufe-hochpunkt-mit-e-funktion-gesucht) [@thispage /XYZ @xpos @ypos null]}
\begin{pfloesungb}{}
\lzb{9.}{}{f$'$ = (2 $-$ x) · $e^{x}$; x = 2; Hochpunkt (2|e²) $\approx$ (2|7,39)}{f$'$(x) = $-e^{x}$ + (3 $-$ x) · $e^{x}$ = (2 $-$ x) · $e^{x}$; $e^x$ > 0 $\Rightarrow$ 2 $-$ x = 0 $\Rightarrow$ x = 2; f(2) = e² $\Rightarrow$ Hochpunkt (2|e²) $\approx$ (2|7,39)}
\end{pfloesungb}
\begin{pfloesungb}{}
\lzb{10.}{}{x = 6; f$'$ wechselt bei 6 das Vorzeichen von + nach $-$; Hochpunkt (6|3)}{$e^{-\tfrac{x}{3} + 2}$ > 0 $\Rightarrow$ f$'$(x) = 0 $\Leftrightarrow$ 2 $-$ x/3 = 0 $\Rightarrow$ x = 6; f$'$ wechselt bei 6 das Vorzeichen von + nach $-$; f(6) = 3 · e$^{0}$ = 3 $\Rightarrow$ Hochpunkt (6|3)}
\end{pfloesungb}
\end{document}